\section{Introducción}

La matriz insumo-producto de Leontief es una herramienta fundamental en el análisis económico que permite estudiar las interacciones entre los diferentes sectores de una economía. Desarrollada por Wassily Leontief en 1936, esta herramienta modela la economía como un sistema de ecuaciones lineales donde la producción de cada sector depende de la demanda final y de los insumos intermedios requeridos por otros sectores.

El departamento de La Guajira, ubicado en la región Caribe de Colombia, presenta una economía caracterizada por la presencia de sectores como la minería, los servicios públicos, la agricultura y el comercio. Comprender las interacciones entre estos sectores es fundamental para la formulación de políticas económicas y la planificación del desarrollo regional.

\subsection{Objetivos}

\textbf{Objetivo general:}

Aplicar la matriz insumo-producto de Leontief a la economía de La Guajira utilizando tres sectores económicos (Agricultura, Tecnología y Servicios) para ilustrar el uso y las bondades de esta herramienta.

\textbf{Objetivos específicos:}

\begin{itemize}
    \item Calcular la matriz de coeficientes técnicos y la matriz inversa de Leontief para los tres sectores seleccionados.
    \item Determinar la producción total requerida dada una demanda final específica.
    \item Analizar el impacto de un aumento en la demanda final de un sector sobre la producción de todos los sectores.
    \item Comparar diferentes escenarios de demanda final y su efecto en la producción requerida.
\end{itemize}

\subsection{Justificación}

La aplicación de la matriz de Leontief a la economía de La Guajira permite comprender cómo los cambios en la demanda de un sector afectan a los demás sectores a través de los encadenamientos productivos. Esta información es valiosa para la toma de decisiones en políticas públicas, la planificación económica y la identificación de sectores estratégicos para el desarrollo regional.

\clearpage
